\documentclass{article}
% \usepackage[notextcomp,nomath]{stix2}
\usepackage{tcolorbox}
\usepackage{amsfonts,amscd,amssymb,amsmath,amsthm,enumerate}
\numberwithin{equation}{section}
\usepackage{array,mathtools,mathrsfs,bm}
\usepackage{threeparttable}
\usepackage{booktabs,caption,longtable,multicol,multirow}
\usepackage{subcaption,lscape}
\usepackage{makecell}
\usepackage{graphicx}
\usepackage{enumerate}
\usepackage{enumitem}
\setlist[enumerate]{noitemsep, topsep=2pt}
\setlist[itemize]{noitemsep, topsep=2pt}
% \usepackage[linesnumbered,boxruled,lined,commentsnumbered]{algorithm2e}
\usepackage[colorlinks,citecolor=purple,linkcolor=blue]{hyperref}
% \usepackage[nameinlink]{cleveref}
\usepackage{cleveref}
% \usepackage[parfill]{parskip}
\usepackage[margin=1.5in]{geometry}
\usepackage[sort&compress,numbers]{natbib}
% \usepackage{authblk}
\usepackage{algorithm,algorithmic}
\usepackage{pifont}
\usepackage{appendix}
\usepackage{multirow}
\usepackage{tabularx}
\usepackage{makecell}
\usepackage{tikz}
\usepackage{pgfplots}
\usetikzlibrary{arrows}
\usetikzlibrary{calc}
\usetikzlibrary{arrows.meta}
\usetikzlibrary{backgrounds}
\usepgfplotslibrary{patchplots}
\usepgfplotslibrary{fillbetween}
% \usepackage[inline]{showlabels}
% \renewcommand{\showlabelfont}{\tiny\color{blue}}
%%%%%%%%%%%%%%%%%%%%%%%%%%%%%%%%%%%%%%%%%%%%%%%%
\usepackage{eqparbox}
\renewcommand\algorithmiccomment[1]{%
  \hfill\texttt{\#}\ \eqparbox{COMMENT}{\texttt{#1}}%
}
%%%%%%%%%%%%%%%%%%%%%%%%%%%%%%%%%%%%%%%%%%%%%%%%

%% tikz
\definecolor{ao(english)}{rgb}{0.0, 0.5, 0.0}
\definecolor{cadmiumgreen}{rgb}{0.0, 0.42, 0.24}
\definecolor{darkpastelgreen}{rgb}{0.01, 0.75, 0.24}
\let\subfigureautorefname=\figurename
\usepackage{adjustbox}
% Define block styles
\tikzset{%
  materia/.style={draw, fill=blue!20, text width=6.0em, text centered, minimum height=1.5em,drop shadow},
  etape/.style={materia, text width=8em, minimum width=10em, minimum height=3em, rounded corners, drop shadow},
  texto/.style={above, text width=6em, text centered},
  linepart/.style={draw, thick, color=black!50, -LaTeX, dashed},
  line/.style={draw, thick, color=black!50, -LaTeX},
  ur/.style={draw, text centered, minimum height=0.01em},
  back group/.style={fill=white!20,rounded corners, draw=black!50, dashed, inner xsep=15pt, inner ysep=10pt},
}
\pgfplotsset{%
  layers/standard/.define layer set={%
      background,axis background,axis grid,axis ticks,axis lines,axis tick labels,pre main,main,axis descriptions,axis foreground%
    }{
      grid style={/pgfplots/on layer=axis grid},%
      tick style={/pgfplots/on layer=axis ticks},%
      axis line style={/pgfplots/on layer=axis lines},%
      label style={/pgfplots/on layer=axis descriptions},%
      legend style={/pgfplots/on layer=axis descriptions},%
      title style={/pgfplots/on layer=axis descriptions},%
      colorbar style={/pgfplots/on layer=axis descriptions},%
      ticklabel style={/pgfplots/on layer=axis tick labels},%
      axis background@ style={/pgfplots/on layer=axis background},%
      3d box foreground style={/pgfplots/on layer=axis foreground},%
    },
}

\newcommand{\etape}[2]{node (p#1) [etape] {#2}}

\newcommand{\transreceptor}[3]{%
  \path [linepart] (#1.east) -- node [above] {\scriptsize #2} (#3);}
\tikzstyle{matheq} = [node distance=8.75cm, text width=21em, minimum width=1cm,
minimum height=2em, text centered,color=black!50]

\renewcommand{\baselinestretch}{1.2}

%
\newcommand{\gkc}{\|\gk\|^{1/2}}
\newcommand{\xk}{x_k}
\newcommand{\Hk}{H_k}
\newcommand{\gk}{g_k}
\newcommand{\dk}{d_k}
\newcommand{\gkn}{g_{k+1}}
\newcommand{\xkn}{x_{k+1}}
\newcommand{\dkn}{d_{k+1}}
\newcommand{\lk}{\lambda_{k+1}}
\newcommand{\cmark}{\ding{51}}%
\newcommand{\xmark}{\ding{55}}%
\newcommand{\arca}{\texttt{CubicReg-A}}

\newtheorem{theorem}{Theorem}[section]
\newtheorem{lemma}[theorem]{Lemma}

\theoremstyle{definition}
\newtheorem{definition}[theorem]{Definition}
\newtheorem{example}[theorem]{Example}
\newtheorem{xca}[theorem]{Exercise}

\theoremstyle{remark}
\newtheorem{remark}[theorem]{Remark}
\newtheorem{corollary}{Corollary}[section]
% \newtheorem{lemma}{Lemma}[section]
\newtheorem{proposition}{Proposition}[section]
\newtheorem{condition}{Condition}[section]
% \newtheorem{definition}{Definition}[section]
% \newtheorem{example}{Example}[section]
\newtheorem{question}{Question}[section]
% \newtheorem{remark}{Remark}[section]
\newtheorem{assumption}{Assumption}[section]
\newtheorem{strategy}{Strategy}[section]


\def\algorithmautorefname{Algorithm}
\def\subsectionautorefname{Section}
\def\definitionautorefname{Definition}
\def\corollaryautorefname{Corollary}
\def\lemmaautorefname{Lemma}
\def\assumptionautorefname{Assumption}
\def\propositionautorefname{Proposition}
\def\tableautorefname{Table}
\def\strategyautorefname{Strategy}
\def\remarkautorefname{Remark}
\def\conditionautorefname{Condition}


\newcommand{\hsodm}{\texttt{HSODM}}
\newcommand{\hsodmarc}{\texttt{HSODMArC}}
\newcommand{\hsodmhvp}{\texttt{HSODM-HVP}}
\newcommand{\drsom}{\texttt{DRSOM}}
\newcommand{\drsomh}{\texttt{DRSOM-H}}
\newcommand{\lbfgs}{\texttt{LBFGS}}
\newcommand{\newtontr}{\texttt{Newton-TR}}
\newcommand{\cg}{\texttt{CG}}
\newcommand{\arc}{\texttt{ARC}}
\newcommand{\atrms}{\texttt{NO-ATR}}
\newcommand{\itrace}{\texttt{TRACE}}
\newcommand{\gd}{\texttt{GD}}
\newcommand{\utr}{\texttt{UTR}}
\newcommand{\utrhvp}{\texttt{UTR-HVP}}
\newcommand{\newtontrst}{\texttt{Newton-TR-STCG}}
\newcommand{\iutr}{\texttt{iUTR(Hessian)}}
\newcommand{\iutrhvp}{\texttt{iUTR}}

\newcommand{\atr}{\texttt{GL-ATR}}

\makeatletter
\renewcommand\paragraph{\@startsection{paragraph}{4}%
  \z@ \z@ {-.5em}%
  {\normalfont\normalsize\bfseries}}%
\makeatother


\usepackage{float}
\usepackage{etoolbox}
\usepackage{cleveref}

% Define subroutine float
\newfloat{subroutine}{htbp}{los}
\floatname{subroutine}{Subroutine}

% Define cref name for subroutine
\crefname{subroutine}{subroutine}{subroutines}
\Crefname{subroutine}{Subroutine}{Subroutines}

\allowdisplaybreaks
\title{Response to the Reviews}
\begin{document}
\maketitle
\section{Reviewer 6LsA}
Thank you for reviewing our paper carefully and evaluating our contribution positively. The question you raise is a great point; we will carefully address your concern as below:
\paragraph{Experimental Designs Or Analyses}

\paragraph{Other Comments Or Suggestions}
\begin{itemize}
    \item Thanks for the suggestion; we will revise our manuscript accordingly.
    \item We are referring to the accelerated CubicReg methods [1], and other accelerated methods generated by plugging second-order oracles into the acceleration framework [2,3,4]. We remark that the framework can only accelerate certain kinds of second-order oracles, with trust-region subproblem notably absent from the list. Please refer to [3] for details. We will polish our text to make it more clear.
    \item We will revise our manuscript as you suggest; please also refer to the answer to your second question below for details. 
    \item We will modify Algorithm 2 to avoid confusion.
    \item Typically, solving a trust-region subproblem involves solving a logarithmic number of linear equations. Therefore, the cost of solving trust-region subproblems is higher than that of linear equations. We will add some references in the manuscript to make the content more clear.
\end{itemize}
\paragraph{Questions For Authors}
\begin{itemize}
    \item We mean in the each iteration of Algorithm 1, all the trust-region subproblems and the linear systems(line 5, alg 1; line 9, alg 2; line 4, alg 3) solved share the same Hessian and gradient, this makes the computational cost lower. In contrast, a series of subproblems with different Hessians and gradients are solved in each iteration of their algorithm.
    \item This is a good point. In fact, as long as $C_1 M \leq \frac{\mu}{\|d\|} \leq C_2 M$ holds for absolution constant $C_1, C_2$, the "sufficient decrease condition" is satisfied. $C_1 = 1,C_2=2$ is a specific choice and a sufficient condition as you said. We can slightly modify the constants in the proof to make them consistent. We will revise our manuscript accordingly.
    \item By inappropriate we mean that the roadmap in Figure 1 fails, and thus it calls for a new choice of subproblem parameters and new algorithm design to accelerate the method.
    \item Yes, the absolute ratio for CubicReg is $\frac{M}{2}$. We plot the normalized ratio in the sense that we compare the current ratio to the average ratio, as can be seen in page 5, line 232.
    \item The reason is that our goal is to find $(\mu,d)$ with $M\leq \frac{\mu}{\|d\|}\leq 2M$ holds. When $\lambda_k^{(j)} >0$, we have $\|d_k^{(j)}\|  = \frac{1}{3\sqrt{M}}\xi^{(j)}$ and $\mu_k^{(j)} := \sigma_k^{(j)}+\lambda_k^{(j)} \geq M\|d_k^{(j)}\|$ from the optimality condition of the subproblem. Therefore, we only need to discuss whether $\frac{\mu_k^{(j)}}{\|d_k^{(j)}\|}\leq 2M$ holds and the corresponding threshold is $\frac{\sqrt{M}}{3}\xi^{(j)}$.
\end{itemize}
[1] Nesterov, Y. Accelerating the cubic regularization of new-
ton’s method on convex problems. Mathematical Pro-
gramming, 112(1):159–181, 2008.
[2] Monteiro, R. D. and Svaiter, B. F. An accelerated hybrid
proximal extragradient method for convex optimization
and its implications to second-order methods. SIAM Jour-
nal on Optimization, 23(2):1092–1125, 2013.
[3] Carmon, Y., Hausler, D., Jambulapati, A., Jin, Y., and Sid-
ford, A. Optimal and adaptive monteiro-svaiter accelera-
tion. Advances in Neural Information Processing Systems,
35:20338–20350, 2022.
[4] Kovalev, D. and Gasnikov, A. The first optimal acceleration
of high-order methods in smooth convex optimization.
Advances in Neural Information Processing Systems, 35:
35339–35351, 2022.
\section{Reviewer hVeX}
Thank you very much for your careful review and positive evaluation of our work. We address your concerns below:
\paragraph{Experimental Designs Or Analyses}
\paragraph{Supplementary Material}
\paragraph{Relation To Broader Scientific Literature}
\paragraph{Essential References Not Discussed}
\paragraph{Other Comments Or Suggestions}
\section{Reviewer 74CJ}
Thank you for reviewing our paper in detail and providing your professional opinions. we will carefully address all the questions you raise and revise our paper accordingly. Below, we will address your questions in detail.
\paragraph{Experimental Designs Or Analyses}
\paragraph{Other Strengths And Weaknesses}
We appreciate your positive feedback on our presentation. The weaknesses you raised are very insightful, we will answer your questions in detail in the "Questions For Authors" section.
\paragraph{Other Comments Or Suggestions}
Sorry for our mistakes! We will fix these typos. Thanks for pointing this out.
\paragraph{Questions For Authors}

\newpage
References:

[1] 

% \bibliographystyle{unsrtnat}
% \bibliography{ref}
\end{document}